%==============================================================
% Curriculum Vitae -- Thevathayarajh Thayananthan
% Built for faculty, postdoctoral, and research scientist applications.
%
% HYPERLINK POLICY (ATS-safe):
%   Every link uses the readable identifier as its anchor text, e.g.
%       \href{https://www.linkedin.com/in/imtheva}{linkedin.com/in/imtheva}
%   NOT \href{...}{LinkedIn}. Parsers read the PDF text layer, so the
%   visible characters must still carry the address. Links are black,
%   so the document looks and prints identically to the plain version.
%
% >>> VERIFY THE MONTHS IN ALL DATE RANGES. <<<
% Months documented from certificates/appointment letters:
%   MAS Intimates internship .......... Sep 2017 -- Feb 2018
%   AHEAD Project Activity 02 ......... from 5 Mar 2019 (one-year term)
%   Robotics short course ............. from 1 Nov 2019
%   Faculty Board ..................... from 13 Sep 2022
%   BGE 121-2 subject coordinator ..... from 20 Sep 2022
%   Employability survey .............. from 9 Nov 2022
% All other months below are best-guess placeholders -- confirm each
% against your transcripts, offer letters, and appointment records.
%==============================================================

\documentclass[10pt, letterpaper]{article}

\usepackage[utf8]{inputenc}
\usepackage[dvipsnames]{xcolor}
\usepackage{etoolbox}
\usepackage{tabularx}
\usepackage{array}
\usepackage{booktabs}
\usepackage{enumitem}
\usepackage{fancyhdr}
\usepackage{lastpage}
\usepackage{titlesec}
\usepackage{amsmath}
\usepackage{calc}
\usepackage{changepage}
\usepackage{paracol}
\usepackage{ifthen}
\usepackage{needspace}
\usepackage{iftex}

\usepackage[
    ignoreheadfoot,
    top=1.7cm,
    bottom=1.7cm,
    left=1.9cm,
    right=1.9cm,
    footskip=0.8cm,
]{geometry}

% All emphasis is now plain bold black. The \textbfr macro is kept only so
% that existing markup keeps working; it no longer applies color.
\newcommand{\textbfr}[1]{\textbf{#1}}

\usepackage[
    pdftitle={Thevathayarajh Thayananthan -- Curriculum Vitae},
    pdfauthor={Thevathayarajh Thayananthan},
    pdfsubject={Curriculum Vitae},
    colorlinks=true, urlcolor=black, linkcolor=black, citecolor=black
]{hyperref}
\usepackage{bookmark}

\ifPDFTeX
    \input{glyphtounicode}
    \pdfgentounicode=1
    \usepackage{lmodern}
\fi

% Shorthand: clickable repository path rendered in typewriter face.
\newcommand{\repo}[1]{\href{https://#1}{\texttt{#1}}}

\AtBeginEnvironment{adjustwidth}{\partopsep0pt}
\setcounter{secnumdepth}{0}
\setlength{\parindent}{0pt}
\setlength{\topskip}{0pt}
\setlength{\columnsep}{0cm}

\pagestyle{fancy}
\fancyhf{}
\fancyfoot[C]{%
  \begin{minipage}[t]{\textwidth}
    \centering
    \begin{tabular*}{\textwidth}{@{\extracolsep{\fill}}lll}
      \small Updated October 2026 &
      \small Thayananthan, Thevathayarajh &
      \small\textit{Page \thepage{} of \pageref*{LastPage}}
    \end{tabular*}
  \end{minipage}%
}
\renewcommand{\headrulewidth}{0pt}
\renewcommand{\footrulewidth}{0pt}

\titleformat{\section}
  {\needspace{4\baselineskip}\bfseries\large}{}{0pt}{}
  [\vspace{1pt}\titlerule]
\titlespacing{\section}{-1pt}{0.30cm}{0.18cm}

\titleformat{\subsection}
  {\needspace{3\baselineskip}\bfseries\normalsize}{}{0pt}{}
\titlespacing{\subsection}{0pt}{0.18cm}{0.08cm}

\renewcommand\labelitemi{$\circ$}

% Give TeX a little slack so long technical terms and URLs never overrun
% the right margin.
\setlength{\emergencystretch}{2em}
\hyphenpenalty=1000

\newenvironment{highlights}{%
    \begin{itemize}[topsep=0.08cm, parsep=0.08cm, partopsep=0pt,
                    itemsep=0pt, leftmargin=0.4cm + 10pt]
}{\end{itemize}}

\newenvironment{onecolentry}{%
    \begin{adjustwidth}{0.2cm}{0.2cm}
}{\end{adjustwidth}}

\newenvironment{twocolentry}[2][]{%
    \onecolentry
    \def\secondColumn{#2}
    \setcolumnwidth{\fill, 4.2cm}
    \begin{paracol}{2}
}{%
    \switchcolumn \raggedleft \secondColumn
    \end{paracol}
    \endonecolentry
}

\newenvironment{pubsJ}{%
  \begin{enumerate}[label={[J\arabic*]}, leftmargin=1.15cm, labelsep=0.2cm,
                    topsep=0.08cm, parsep=0.08cm, partopsep=0pt, itemsep=0pt]
}{\end{enumerate}}

\newenvironment{pubsB}{%
  \begin{enumerate}[label={[B\arabic*]}, leftmargin=1.15cm, labelsep=0.2cm,
                    topsep=0.08cm, parsep=0.08cm, partopsep=0pt, itemsep=0pt]
}{\end{enumerate}}

\newenvironment{pubsC}{%
  \begin{enumerate}[label={[C\arabic*]}, leftmargin=1.15cm, labelsep=0.2cm,
                    topsep=0.08cm, parsep=0.08cm, partopsep=0pt, itemsep=0pt]
}{\end{enumerate}}

\newenvironment{pubsS}{%
  \begin{enumerate}[label={[S\arabic*]}, leftmargin=1.15cm, labelsep=0.2cm,
                    topsep=0.08cm, parsep=0.08cm, partopsep=0pt, itemsep=0pt]
}{\end{enumerate}}

\newenvironment{pubsU}{%
  \begin{enumerate}[label={[U\arabic*]}, leftmargin=1.15cm, labelsep=0.2cm,
                    topsep=0.08cm, parsep=0.08cm, partopsep=0pt, itemsep=0pt]
}{\end{enumerate}}

\newenvironment{pres}{%
  \begin{enumerate}[label={[P\arabic*]}, leftmargin=1.15cm, labelsep=0.2cm,
                    topsep=0.08cm, parsep=0.08cm, partopsep=0pt, itemsep=0pt]
}{\end{enumerate}}

\let\hrefWithoutArrow\href

%==============================================================
\begin{document}

%--------------------------------------------------------------
% HEADER -- name, then contact details stacked and centered.
%--------------------------------------------------------------
\begin{center}
    {\fontsize{16pt}{18pt}\selectfont\bfseries Thevathayarajh Thayananthan}

    \vspace{0.18cm}

    Riverbend Research Lab North, Room 001E, 110 Riverbend Road, Athens, GA 30605 \\
    +1 662 370 6844 \\
    \href{mailto:theva@uga.edu}{theva@uga.edu} \\
    \href{https://imtheva.github.io}{imtheva.github.io} \textbar{}
    \href{https://www.linkedin.com/in/imtheva}{linkedin.com/in/imtheva} \textbar{}
    \href{https://orcid.org/0009-0007-2576-4348}{ORCID 0009-0007-2576-4348} \textbar{}
    \href{https://scholar.google.com/citations?user=MpKhKEUAAAAJ}{Google Scholar}
\end{center}

\vspace{0.25cm}

%--------------------------------------------------------------
\section{RESEARCH FOCUS}
\vspace{0.06cm}
\begin{onecolentry}
    Research focuses on AI-enabled autonomous systems for agriculture, integrating
    computer vision, multimodal sensing, 3D localization, autonomous field navigation,
    and robotic manipulation for selective harvesting and precision crop operations.
    Develops complete perception-to-action robotic systems, from ROS and Gazebo
    simulation to integrated hardware and field deployment, to advance reliable,
    practical, and affordable agricultural automation.
\end{onecolentry}
\vspace{0.18cm}

%--------------------------------------------------------------
\section{EDUCATION}
\vspace{0.08cm}

\begin{twocolentry}{Aug 2025 -- Present}
    \textbf{Ph.D., Agricultural Engineering} \\
    University of Georgia, Athens, GA, USA \\
    Advanced to candidacy August 2026; \textbf{Expected graduation May 2027} \\
    Dissertation Proposal: \emph{Design, Development, Integration and Field Evaluation of a
    Robotic Cotton Picker}
    \vspace{0.12cm}
\end{twocolentry}

\begin{twocolentry}{Jan 2023 -- Jul 2025}
    \textbf{Doctoral coursework and research, Biosystems Engineering} \\
    Mississippi State University, Mississippi State, MS, USA \\
    Program transferred to the University of Georgia with advisor
    \vspace{0.12cm}
\end{twocolentry}

\begin{twocolentry}{Jan 2019 -- Jul 2022}
    \textbf{M.Sc., Electrical and Electronics Engineering} \\
    University of Peradeniya, Sri Lanka \\
    Project: IoT-based medical instrumentation for remote patient monitoring
    \vspace{0.12cm}
\end{twocolentry}

\begin{twocolentry}{Apr 2014 -- Jan 2018}
    \textbf{B.Tech.\ (Hons.), Science and Technology --- Mechatronics} \\
    Uva Wellassa University of Sri Lanka, Badulla, Sri Lanka \\
    Dissertation: \emph{Development of a monitoring device for the fermentation stage of
    black tea manufacturing}
    \vspace{0.08cm}
\end{twocolentry}

%--------------------------------------------------------------
\section{RESEARCH APPOINTMENTS}
\vspace{0.08cm}

\begin{twocolentry}{Aug 2025 -- Present}
    \textbf{Graduate Research Assistant} \\
    Sensing and Automation in Agri-Systems (SAAS) Lab, School of Environmental,
    Civil, Agricultural \& Mechanical Engineering, University of Georgia
    \begin{highlights}
        \item Develop perception-driven manipulator control in ROS and Python for a
              dual-side robotic cotton picker. Completed field testing with one
              ZED~2i stereo camera, a UR5e manipulator, and a custom end-effector,
              and developed a two-camera control pipeline for dual-side harvesting.

        \item Benchmark 33 YOLO variants (v8--v13) and several segmentation models on
              a 1{,}008-image, four-class annotated field dataset, comparing mAP and
              inference speed on matched test splits.

        \item Developed \textbf{CottonID}, a full-stack web application using Next.js,
              FastAPI, and SQLite to generate boll counts, yield estimates, and
              defoliation-timing recommendations. Currently improving the underlying
              perception models; deployed through Cloudflare Tunnel for
              demonstrations.

        \item Manage equipment, operations, data infrastructure, and safety within the
              Sensing \& Automation in Agri-Systems (SAAS) Lab.
    \end{highlights}
    \vspace{0.12cm}
\end{twocolentry}

\begin{twocolentry}{Jan 2023 -- Jul 2025}
    \textbf{Graduate Research Assistant} \\
    Department of Agricultural and Biological Engineering,
    Mississippi State University
    \begin{highlights}
        \item Designed and open-sourced \textbf{CottonSim}, a ROS/Gazebo simulation
              environment integrating a virtual cotton farm, boll perception, and
              vision-guided navigation. The Cotton-Eye module (YOLOv8n-seg) reached
              85.2\% mAP, 88.9\% recall, and 93.0\% precision; navigation achieved
              100\% GPS-guided and 96.7\% map-based traversal success within
              0.25~m~[J1].

        \item Developed an in-field multi-ripeness blackberry-detection pipeline for
              soft robotic harvesting from 1{,}086 images; benchmarking nine YOLO
              architectures, where YOLOv7-base produced 91.1\% mAP overall, 92.4\%
              for ripe berries, and 12.6~ms inference per
              $1024 \times 1024$ image~[J2].

        \item Developed a semantic-segmentation pipeline for automating bladeless
              catfish cutting, benchmarking five architectures across five classes on
              396 images; SegFormer-B5 reached 89.2\% mIoU and 94.6\% mean pixel
              accuracy~[C3].

        \item Named as an inventor on a university technology disclosure for an
              autonomous vision-guided cotton-harvesting system.
    \end{highlights}
    \vspace{0.12cm}
\end{twocolentry}

\begin{twocolentry}{Sep 2017 -- Feb 2018}
    \textbf{Research and Development Intern} \\
    Research \& Innovation Department, MAS Intimates
    (Unichela Pvt.\ Ltd., Vidiyal Division), Sri Lanka
    \begin{highlights}
        \item Contributed to seven industrial-automation projects in garment
              manufacturing, including PLC-based finished-goods scanning and garment
              folding, an Arduino-based chiller-shutdown SMS alert system,
              pedal-free and automated sewing, and waistband cutting with mark
              elimination.
    \end{highlights}
    \vspace{0.08cm}
\end{twocolentry}

%--------------------------------------------------------------
\section{ACADEMIC APPOINTMENTS}
\vspace{0.08cm}

\begin{twocolentry}{Mar 2019 -- Dec 2022}
    \textbf{Lecturer (Full-time)} \\
    Department of Science and Technology, Faculty of Applied Sciences,
    Uva Wellassa University of Sri Lanka
    \begin{highlights}
        \item Instructor of record for eight distinct undergraduate courses totaling 18 credit hours across four academic years (see Teaching)
              % TODO: typical class size, e.g. "class sizes of 30-60 students".

        \item Designed curricula for seven new undergraduate courses for the
              B.Sc.\ in Mechatronics curriculum revision, including Computer
              Vision and IoT Development.

        \item Supervised 15 undergraduate final-year research projects, twelve of
              which produced peer-reviewed conference abstracts/papers~[U1]--[U12].
    \end{highlights}
    \vspace{0.12cm}
\end{twocolentry}

\begin{twocolentry}{Mar 2018 -- Mar 2019}
    \textbf{Demonstrator (Temporary)} \\
    Department of Science and Technology, Faculty of Applied Sciences,
    Uva Wellassa University of Sri Lanka
    \begin{highlights}
        \item Led undergraduate laboratory and practical sessions in electronics,
              instrumentation, and workshop practice, supervising groups of
              students in the lab class and grading laboratory reports.

        \item Maintained teaching laboratory equipment and assisted faculty with
              demonstrations, practical assessments, and student project
              supervision.
    \end{highlights}
    \vspace{0.08cm}
\end{twocolentry}

%--------------------------------------------------------------
\section{PUBLICATIONS}
\vspace{0.04cm}
\begin{onecolentry}
    {\footnotesize Name in \textbf{bold}; $^{*}$corresponding author.
    Earlier work appears under the name \textbf{Thevathayarajh, T.}
    Citations: 41 \textbar{} h-index: 4 (Google Scholar).\par}
\end{onecolentry}
\vspace{0.12cm}

\subsection{Peer-Reviewed Journal Articles}
\begin{onecolentry}
\begin{pubsJ}
    \item \textbfr{Thayananthan, T.}, Zhang, X.$^{*}$, Huang, Y., Chen, J.,
    Wijewardane, N.K., Martins, V.S., Chesser, G.D., and Goodin, C.T. (2025).
    CottonSim: A vision-guided autonomous robotic system for cotton harvesting in
    Gazebo simulation. \emph{Computers and Electronics in Agriculture}, 239, 110963.
    \url{https://doi.org/10.1016/j.compag.2025.110963}

    \item \textbfr{Thayananthan, T.}, Zhang, X.$^{*}$, Harjono, J., Huang, Y., Liu, W.,
    McWhirt, A.L., Threlfall, R.T., and Chen, Y. (2025). In-field multi-ripeness
    blackberry detection for soft robotic harvesting. \emph{Journal of the ASABE},
    68(6), 1073--1089. \url{https://doi.org/10.13031/ja.16300}
\end{pubsJ}
\end{onecolentry}
\vspace{0.12cm}

\subsection{Book Chapters}
\begin{onecolentry}
\begin{pubsB}
    \item Gunaratnam, A., \textbfr{Thayananthan, T.}, Thangathurai, K., and
    Abhiram, B. (2024). Computer vision in livestock management and production.
    In \emph{Engineering Applications in Livestock Production}, pp.~93--128.
    Academic Press. \url{https://doi.org/10.1016/B978-0-323-98385-3.00002-5}
\end{pubsB}
\end{onecolentry}
\vspace{0.12cm}

\subsection{Peer-Reviewed Conference Proceedings}
\begin{onecolentry}
\begin{pubsC}
    \item \textbfr{Thayananthan, T.} and Zhang, X.$^{*}$ (2025).
    Perception-enabled manipulator control for a robotic cotton picker with
    dual-side harvesting capability. \emph{IFAC-PapersOnLine}, 59(23), 332--337.
    (8th IFAC Conference on Sensing, Control and Automation Technologies for
    Agriculture, AGRICONTROL 2025, UC Davis, CA, USA.)
    \url{https://doi.org/10.1016/j.ifacol.2025.11.809}

    \item Vimalambikaipakan, G., Amarasinghe, C., Rajapaksha, T.,
    \textbfr{Thayananthan, T.}, and Mariyathas, J. (2024).
    Traffic sign recognition and auditory alert system for Sri Lankan drivers using
    deep learning. In \emph{Proc.\ International Research Conference on Smart
    Computing and Systems Engineering (SCSE)}, 7, 1--5. IEEE.
    \url{https://doi.org/10.1109/SCSE61872.2024.10550615}

    \item \textbfr{Thayananthan, T.}, Zhang, X.$^{*}$, Liu, W., Yao, T., Huang, Y.,
    Wijewardane, N.K., and Lu, Y. (2023). Automating catfish cutting process using
    deep learning-based semantic segmentation. In \emph{Sensing for Agriculture and
    Food Quality and Safety XV}, 12545, 103--116. SPIE.
    \url{https://doi.org/10.1117/12.2663370}

    \item Zhang, X.$^{*}$, \textbfr{Thayananthan, T.}, Usman, M., Liu, W., and Chen, Y.
    (2023). Multi-ripeness level blackberry detection using YOLOv7 for soft robotic
    harvesting. In \emph{Autonomous Air and Ground Sensing Systems for Agricultural
    Optimization and Phenotyping VIII}, 12539, 85--96. SPIE.
    \textbf{Best Paper Award.} \url{https://doi.org/10.1117/12.2663367}

    \item Yapa, A., Ratnayake, K., Premarathna, C.P., and \textbfr{Thayananthan, T.}
    (2023). Leveraging virtual reality for robot manipulator education. In
    \emph{Proc.\ Moratuwa Engineering Research Conference (MERCon)}, 544--549. IEEE.
    \url{https://doi.org/10.1109/MERCon60487.2023.10355455}

    \item \textbfr{Thayananthan, T.}, Wanniarachchi, D.D.C., and Kumari, K.W.S.N.
    (2019). A digital device to monitor the colour changes during the fermentation
    stage of black tea processing. In \emph{Proc.\ International Conference on Recent
    Advances in Computer Science and Engineering (IC-RACE)}, 75--79.
\end{pubsC}
\end{onecolentry}
\vspace{0.12cm}

\subsection{Preprints and Manuscripts in Preparation or Review}
\begin{onecolentry}
\begin{pubsS}
    \item \textbfr{Thayananthan, T.}, Zhang, X.$^{*}$, Laddusinghe Badu, I.,
    Harjono, J., Rains, G.C., Li, B., Bastos, L.M., Wijewardane, N.K., and
    Martins, V.S. (2026). Selective cotton boll localization for robotic harvesting:
    evaluation of deep learning vision models under field conditions.
    \emph{arXiv preprint} arXiv:2609.19592. Under preparation for journal
    submission. \url{https://arxiv.org/abs/2609.19592}

    \item \textbfr{Thayananthan, T.} and Zhang, X. Toward fully autonomous cotton
    harvesting robots: a systematic review of perception, navigation, path planning,
    and harvesting mechanisms. \emph{Under internal review} (2026).

    \item \textbfr{Thayananthan, T.} and Zhang, X. Field integration and evaluation of
    an autonomous robotic cotton picker. \emph{In preparation} (2026).
\end{pubsS}
\end{onecolentry}
\vspace{0.12cm}

\subsection{Supervised Undergraduate Research --- Conference Abstracts/Papers}
\begin{onecolentry}
\begin{pubsU}
    \item Ahamed, M.M.R., Amarasinghe, C., \textbfr{Thayananthan, T.},
    De Silva, D.S.P.P.G., De Zoysa, K., and De Silva, D.P.N. (2023). Development of
    automated image capturing for zebrafish embryo toxicity model.
    \emph{Proc.\ International Research Conference, Uva Wellassa University}.

    \item Geerthana, V., \textbfr{Thayananthan, T.}, Amarasinghe, A.R.P.C.C.J.,
    Janotheepan, M., and Lakmali, R.M.T. (2023). Deep learning-based traffic sign
    recognition and auditory alert system for Sri Lankan drivers.
    \emph{Proc.\ International Research Conference, Uva Wellassa University}.

    \item Liyadipita, L.A.M.G.S.K., Amarasingha, A.R.P.C.C.J.,
    \textbfr{Thayananthan, T.}, and Bandara, T.U.K.S. (2023). IoT-based patient
    monitoring system for Sri Lanka. \emph{Proc.\ International Research Conference,
    Uva Wellassa University}.

    \item Jayawardhana, D.G.H., \textbfr{Thayananthan, T.}, Ekanayake, R.M.T.C.B.,
    and Senanayake, D.D.B. (2022). Automatic billing system for a Sri Lankan pastry
    shop. \emph{Proc.\ International Research Conference, Uva Wellassa University}.

    \item Alwis, D.M.G.P., \textbfr{Thayananthan, T.}, Ekanayake, R.M.T.C.B., and
    Senanayake, D.D.B. (2022). Motorbike assistance tool using image processing
    techniques. \emph{Proc.\ International Research Conference, Uva Wellassa
    University}.

    \item Sampath, W.P.A.C., \textbfr{Thayananthan, T.}, Prematilake, M., and
    Ekanayake, R.M.T.C.B. (2022). Development of an IoT-based automated colony
    counter. \emph{Proc.\ International Research Conference, Uva Wellassa
    University}.

    \item Arawa, A.M.P.R.B., \textbfr{Thayananthan, T.}, Mehendran, Y., and
    Senanayake, D.B.B. (2022). Recognition of Sinhala machine-printed text for
    postal address interpretation and postal automation.
    \emph{Proc.\ International Research Conference, Uva Wellassa University}.

    \item Dissanayaka, D.M.K.C., Vidanapathirana, A.C., and
    \textbfr{Thayananthan, T.} (2022). Remote monitoring and controlling of an
    automobile system using NodeMCU. \emph{Proc.\ International Research Conference,
    Uva Wellassa University}.

    \item Gayashan, S.A.V., Vidanapathirana, A.C., Ekanayake, R.M.T.C.B., and
    \textbfr{Thayananthan, T.} (2022). Development of a wind turbine system for
    electric and hybrid carts. \emph{Proc.\ International Research Conference,
    Uva Wellassa University}.

    \item Srikaran, K. and \textbfr{Thayananthan, T.} (2021). Development of an
    IoT-based tree-cutting warning system in forests. \emph{Proc.\ 3rd South Asia
    Conference on Interdisciplinary Research}, 14.

    \item Danijal, T.J. and \textbfr{Thayananthan, T.} (2021). Development of a
    tracking over-speed system using IoT technology for vehicles.
    \emph{Proc.\ International Research Conference, Uva Wellassa University}, 139.

    \item Senanayake, D.D.B., Kumari, K.W.S.N., and \textbfr{Thayananthan, T.}
    (2020). Active and passive safety system for differently abled people and adults.
    \emph{Proc.\ International Research Conference, Uva Wellassa University}, 195.

    \item \textbfr{Thayananthan, T.}, Amarasinghe, A.R.P.C.C.J., and
    Ekanayake, R.M.T.C.B. (2019). A digital device to measure distance on multiple
    surfaces. \emph{Proc.\ International Research Conference, Uva Wellassa
    University}, 504.
\end{pubsU}
\end{onecolentry}
\vspace{0.18cm}

%--------------------------------------------------------------
\section{CONFERENCE PRESENTATIONS}
\vspace{0.08cm}
\begin{onecolentry}
\begin{pres}
    \item \textbfr{Thayananthan, T.} (2026, July). Develop and evaluate a simplified
    end-effector for robotic selective cotton harvesting. ASABE Annual International
    Meeting, Indianapolis, IN, USA.

    \item \textbfr{Thayananthan, T.} (2026, March). Harvesting cotton before it's too
    late: a smarter robotic approach. UGA Institute for Integrative Precision
    Agriculture Spring Retreat, Griffin, GA, USA.
    \textbf{Third Place, Three Minute Thesis oral presentation.}

    \item \textbfr{Thayananthan, T.} (2025, August). Perception-enabled manipulator
    control for a robotic cotton picker with dual-side harvesting capability.
    8th IFAC Conference on Sensing, Control and Automation Technologies for
    Agriculture (AGRICONTROL 2025), Davis, CA, USA.

    \item \textbfr{Thayananthan, T.} (2025, April). Cotton-SAM: AI-driven cotton boll
    localization for autonomous picking using YOLO and SAM integration.
    AI in Agriculture and Natural Resources Conference, Starkville, MS, USA.

    \item \textbfr{Thayananthan, T.} (2024, October). Cotton boll localization system
    to enable autonomous cotton picking using YOLO and SAM. Graduate Student Research
    Symposium, Mississippi State University, Mississippi State, MS, USA.

    \item \textbfr{Thayananthan, T.} (2024, July). Development of a cotton boll
    detection system to enhance autonomous picking using YOLOv8 and SAM.
    ASABE Annual International Meeting, Anaheim, CA, USA.

    \item \textbfr{Thayananthan, T.} (2024, July). Computer vision-enabled autonomous
    robotic navigation system for cotton farms in a Gazebo simulation environment.
    ASABE Annual International Meeting, Anaheim, CA, USA.

    \item \textbfr{Thayananthan, T.} (2024, March). Which is better cotton picker?
    Three Minute Thesis presentation, Mississippi Agriculture Consortium,
    Mississippi State University, Starkville, MS, USA.

    \item \textbfr{Thayananthan, T.} (2023, July). Catfish cutting automation through
    computer vision and deep learning. Mississippi Academy of Sciences 5th Summer
    Science \& Engineering Symposium, Mississippi State, MS, USA.

    \item \textbfr{Thayananthan, T.} (2023, May). Automating catfish cutting process
    using deep learning-based semantic segmentation. \emph{Sensing for Agriculture
    and Food Quality and Safety XV} session, SPIE Defense + Commercial Sensing,
    Orlando, FL, USA.

    \item \textbfr{Thayananthan, T.} (2023, May). Multi-ripeness level blackberry
    detection using YOLOv7 for soft robotic harvesting. \emph{Autonomous Air and
    Ground Sensing Systems for Agricultural Optimization and Phenotyping VIII}
    session, SPIE Defense + Commercial Sensing, Orlando, FL, USA.
    \textbf{Best Paper Award.}

    \item \textbfr{Thayananthan, T.} (2018, February). Development of a monitoring
    device for the fermentation stage of black tea manufacturing.
    \emph{Mechanical Engineering and Mechatronics} session, 2nd International
    Research Symposium, Uva Wellassa University, Badulla, Sri Lanka.
\end{pres}
\end{onecolentry}
\vspace{0.18cm}

%--------------------------------------------------------------
\section{HONORS AND AWARDS}
\vspace{0.08cm}
\begin{onecolentry}
    \begin{highlights}
        \item \textbf{Third Place, Three Minute Thesis (3MT) Competition} ---
              Institute for Integrative Precision Agriculture Retreat, University of
              Georgia (2026). Includes a \$500 conference travel award.
        \item \textbf{Hall of Fame Awardee}, College of Engineering ---
              The Graduate School, Mississippi State University (2025)
        \item \textbf{SEC Emerging Scholar} --- Selected to represent Mississippi
              State University in the Southeastern Conference Emerging Scholars
              Program (2023--2024)
        \item \textbf{Best Paper Award} --- Autonomous Air and Ground Sensing Systems
              for Agricultural Optimization and Phenotyping VIII, SPIE (2023)
    \end{highlights}
\end{onecolentry}
\vspace{0.18cm}

%--------------------------------------------------------------
\section{FUNDING AND INTELLECTUAL PROPERTY}
\vspace{0.08cm}
\begin{onecolentry}
    \textbf{Intellectual Property}
    \begin{highlights}
        \item \textbf{Technology Disclosure} --- ``Design and development of an
              autonomous vision-guided cotton picker for selective and targeted
              robotic picking.'' Mississippi State University Office of Technology
              Management (2025). Co-inventor with Dr.\ Xin Zhang.
              \\ \href{https://msstate-innovations.technologypublisher.com/technology/58899}{\texttt{msstate-innovations.technologypublisher.com/technology/58899}}
    \end{highlights}
    \vspace{0.10cm}
    \textbf{Competitive Travel Awards}
    \begin{highlights}
        \item CENGR Graduate Student Travel Grant, College of Engineering,
              University of Georgia (2026) --- \$735
        \item IIPA 3MT Conference Travel Award, University of Georgia (2026) --- \$500
        \item TAGGS Travel Scholarship, The Graduate School, Mississippi State
              University (2024) --- \$700
        \item AG2PI Travel Award, Agricultural Genome to Phenome Initiative
              (2024) --- \$1{,}500
    \end{highlights}
    \vspace{0.10cm}
    \textbf{Proposal Experience}
    \begin{highlights}
        \item Southern SARE Graduate Student Grant (2024) --- \$22{,}000;
              submitted, not funded.
    \end{highlights}
\end{onecolentry}
\vspace{0.18cm}

%--------------------------------------------------------------
\section{TEACHING EXPERIENCE}
\vspace{0.08cm}

\subsection{Instructor of Record --- Uva Wellassa University of Sri Lanka (Mar 2019 -- Dec 2022)}
\begin{onecolentry}
Sole instructor for eight undergraduate courses (18 credit hours total): lectures,
laboratory design, assessment, and grading.
\begin{highlights}
    \item \textbf{SCT 364-2 Embedded Systems} (2 credits) --- 2019, 2020, 2021, 2022
    \item \textbf{SCT 462-2 Intelligent Systems} (2 credits) --- 2019, 2020, 2021, 2022
    \item \textbf{SCT 319-2 Maintenance, Production, and Project Management}
          (2 credits) --- 2021, 2022
    \item \textbf{SCT 363-2 Electrical Systems} (2 credits) --- 2020, 2021
    \item \textbf{SCT 408-2 Industrial Electronics} (2 credits) --- 2020, 2021
    \item \textbf{SCT 362-3 Digital and Analog Electronics} (3 credits) --- 2020
    \item \textbf{SCT 373-2 Electric Power and Machines} (2 credits) --- 2019
    \item \textbf{SCT 142-3 Engineering Workshop} (3 credits) --- 2019
\end{highlights}
\end{onecolentry}
\vspace{0.14cm}

\subsection{Course Development}
\begin{onecolentry}
    \textbf{BSc Mechatronics curriculum revision}, Department of Science and
    Technology, Uva Wellassa University (2019). Invited by the department to design
    seven new course modules and their lesson plans:
    \begin{highlights}
        \item MEC 361-2 Computer Vision
        \item MEC 331-2 IoT Development
        \item MEC 232-2 Embedded Systems
        \item MEC 233-2 Embedded Systems Design
        \item MEC 362-2 Transducers and Instrumentation
        \item MEC 111-2 Computer Programming I
        \item MEC 212-2 Computer Programming II
    \end{highlights}
\end{onecolentry}
\vspace{0.14cm}

\subsection{Course Coordination and Laboratory Instruction}
\begin{onecolentry}
    \begin{highlights}
        \item \textbf{Graduate Teaching Assistant}, ABE 4423/6423 Bioinstrumentation~II,
              Mississippi State University (Spring 2024). Laboratory instructor with
              Dr.\ Xin Zhang. Certified Teaching Assistant, Level~2/3.
        \item \textbf{Subject Coordinator}, BGE 121-2 Ethics and Law Basics,
              Uva Wellassa University (2022). Cross-faculty course delivered to three
              faculties; coordinated visiting resource persons, assessment, and
              results submission.
        \item \textbf{Coordinator}, Short Course in Robotics and Arduino Programming,
              Centre for Open and Distance Learning, Uva Wellassa University
              (2019--2022). Appointed by the Vice Chancellor; responsible for
              curriculum, timetabling, certification, and course budget.
    \end{highlights}
\end{onecolentry}
\vspace{0.14cm}

\subsection{Courses Prepared to Teach}
\begin{onecolentry}
    Computer Vision for Agricultural and Biological Systems; Artificial Intelligence
    for Agricultural and Biological Engineering; Machine Learning Applications in
    Agriculture; Agricultural Automation and Precision Agriculture.
\end{onecolentry}
\vspace{0.18cm}

%--------------------------------------------------------------
\section{MENTORING AND ADVISING}
\vspace{0.08cm}
\begin{onecolentry}
    \textbf{Undergraduate and Visiting Researchers} --- six students mentored, Sensing
    and Automation in Agri-Systems (SAAS) Lab (2024 -- Present)
    \begin{highlights}
        \item \textbf{Isadora Goncalves} (Aug 2026 -- present) --- School of Electrical
              \& Computer Engineering, University of Georgia. End-effector task design
              for a robotics competition platform.
        \item \textbf{Jap Ja Mun Pan} (May -- Jul 2026) --- Electrical \& Computer
              Engineering, Trevecca Nazarene University; visiting researcher through
              the Summer Undergraduate Research Opportunities in Engineering (SUROE)
              program, University of Georgia. Actuation design for a plant surgical
              robot.
        \item \textbf{Madison Mateyko} (Nov -- Dec 2025) --- School of Environmental,
              Civil, Agricultural \& Mechanical Engineering, University of Georgia.
              Computer vision for the robotic cotton picker.
        \item \textbf{Riley Gebke} (Mar -- May 2025) --- Department of Agricultural \&
              Biological Engineering, Mississippi State University. Computer vision
              for the robotic cotton picker.
        \item \textbf{Mary J. Echelberry} (Oct -- Dec 2024) --- Department of
              Agricultural \& Biological Engineering, Mississippi State University.
              Computer vision for the robotic cotton picker.
        \item \textbf{Kelly Bilington} (Apr 2024) --- Harper Adams University, United
              Kingdom; graduate student worker through the Mississippi State
              University--Harper Adams University Student Exchange Program, hosted by
              the Agricultural Autonomy Institute. Development of a UR5e manipulator
              base platform for integration with a Husky A200 mobile robot.
    \end{highlights}
    \vspace{0.12cm}

    \textbf{Undergraduate Final-Year Research Projects} --- Uva Wellassa University
    (Mar 2019 -- Dec 2022). Supervised 15 projects; twelve produced peer-reviewed
    conference abstracts~[U1]--[U12]. Topics spanned computer vision, deep learning,
    IoT systems, and mechatronic design, including traffic sign recognition,
    zebrafish embryo toxicity imaging, automated microbial colony counting, Sinhala
    postal-address OCR, and vehicle safety systems.
    \vspace{0.12cm}

    \textbf{Student Support Roles} --- Uva Wellassa University
    \begin{highlights}
        \item \textbf{Student Counselor}, Faculty of Applied Sciences (2022--2023).
              Completed a 40-hour Psychological Counselling Program.
        \item \textbf{Academic Sub-Warden}, University residence halls (2019--2023).
              Residential academic support and student welfare.
        \item \textbf{Student Mentor}, Faculty of Applied Sciences (2019--2023).
              Academic guidance for undergraduate students.
    \end{highlights}
\end{onecolentry}
\vspace{0.18cm}

%--------------------------------------------------------------
\section{PROFESSIONAL SERVICE}
\vspace{0.08cm}
\begin{onecolentry}
    \textbf{Editorial Board Service}
    \begin{highlights}
        \item \textbf{Associate Editor, Computer Science and Engineering} ---
              \emph{Endeavors: Mississippi State Undergraduate Research Journal}
              (ISSN 3071-012X), Mississippi State University (2024--2025).
              Inaugural editorial board of MSU's first peer-reviewed undergraduate
              research journal; conducted double-blind peer review of engineering
              and computer science submissions for Volume~1.
              \\ \href{https://scholarsjunction.msstate.edu/endeavors/}{\texttt{scholarsjunction.msstate.edu/endeavors}}
    \end{highlights}
    \vspace{0.10cm}
    \textbf{Journal Peer Review} --- 14 completed reviews across seven journals
    \begin{highlights}
        \item \emph{Computers and Electronics in Agriculture}, Elsevier
              (6 reviews; 2024--present)
        \item \emph{The International Journal of Robotics Research}, SAGE
              (1 review; 2026)
        \item \emph{Journal of Food Composition and Analysis}, Elsevier (2 reviews; 2026)
        \item \emph{Frontiers in Plant Science}, Frontiers (2 reviews; 2026)
        \item \emph{Food and Energy Security}, Wiley (1 review; 2026)
        \item \emph{Food and Humanity}, Elsevier (1 review; 2026)
        \item \emph{Journal of Atmospheric and Solar-Terrestrial Physics}, Elsevier
              (1 review; 2026)
    \end{highlights}
    \vspace{0.10cm}
    \textbf{Conference Service}
    \begin{highlights}
        \item Track Coordinator, Engineering and Technology track, International
              Research Conference, Uva Wellassa University (2022)
        \item Reviewer, International Research Conference, Uva Wellassa University
              (2021--2022)
        \item Reviewer, 2nd International Conference on Advancements in Computing,
              Sri Lanka Institute of Information Technology (2021)
        \item Session Panel Member, International Research Conference,
              Uva Wellassa University (2020)
        \item Organizing Committee Member, International Research Conference,
              Uva Wellassa University (2019, 2020)
    \end{highlights}
\end{onecolentry}
\vspace{0.18cm}

%--------------------------------------------------------------
\section{INSTITUTIONAL SERVICE AND GOVERNANCE}
\vspace{0.08cm}
\begin{onecolentry}
    \begin{highlights}
        \item \textbf{Lab Manager}, SAAS Lab --- Mississippi State University
              (2024--2025); University of Georgia (2025--present)
        \item \textbf{Member, Faculty Board}, Faculty of Applied Sciences,
              Uva Wellassa University (2022--2023). Appointed by the Dean under
              S.48(1A)(c) of the Universities Act No.~16 of 1978. The Faculty Board is
              the governing academic body for curricula, examinations, and policy.
              Three-year term, curtailed on departure for doctoral study.
        \item \textbf{Member}, Program Review and Self-Evaluation Report (SER) groups
              3, 4, and 8, Uva Wellassa University (2022--2023)
        \item \textbf{Member}, Graduate Employability Survey team, AHEAD ELTA-ELSE
              Faculty Development Plan, Uva Wellassa University (2022--2023).
              Appointed by the Dean; data informed the final AHEAD project report.
        \item \textbf{Faculty Advisor}, IEEE Robotics and Automation Society Student
              Chapter, Uva Wellassa University (2021--2023)
        \item \textbf{Panel Member}, Stakeholder Meeting for the BSc Mechatronics
              Curriculum Revision, Uva Wellassa University (2020)
        \item \textbf{Co-author}, Survey of Postgraduate Admission Pathways followed by
              UWU graduates (2020)
        \item \textbf{Member (Activity 02)}, AHEAD Project --- Revising the Curricula
              of Existing Degree Programmes to Meet Local and International Demands,
              Uva Wellassa University (2019--2020)
        \item \textbf{Organizing Committee Member}, Convocation, Uva Wellassa
              University (2019)
    \end{highlights}
\end{onecolentry}
\vspace{0.18cm}

%--------------------------------------------------------------
\section{OUTREACH AND COMMUNITY ENGAGEMENT}
\vspace{0.08cm}
\begin{onecolentry}
    \begin{highlights}
        \item \textbf{Secretary}, Softball Cricket Club, Mississippi State University
              (2024--2025)
        \item \textbf{Social Secretary}, Softball Cricket Club, Mississippi State
              University (2023--2024)
        \item \textbf{Webmaster}, Sri Lankan Association, Mississippi State University
              (2023--2024)
        \item \textbf{Workshop Instructor}, Engineering Machinery Workshop for School
              Students, Uva Wellassa University (2018)
        \item \textbf{Organizing Committee}, UWU Robot Battles, Department of Science
              \& Technology, Uva Wellassa University (2016)
        \item \textbf{Volunteer Mathematics Instructor}, G.C.E. Ordinary Level (O/L)
              Seminar Series, Uva Wellassa University, Sri Lanka (2014--2017).
              Conducted mathematics seminars for students at resource-constrained
              schools with limited access to supplemental educational support, with
              the goal of improving examination preparedness, pass rates, and student
              interest in mathematics.
    \end{highlights}
\end{onecolentry}
\vspace{0.18cm}

%--------------------------------------------------------------
\section{PROFESSIONAL DEVELOPMENT}
\vspace{0.08cm}
\begin{onecolentry}
    \begin{highlights}
        \item \textbf{Peer Review Certification} --- Elsevier Researcher Academy
              (2026): Fundamentals of Peer Review; Certified Peer Reviewer Course;
              Becoming a Peer Reviewer (14 hours total)
        \item \textbf{Preparing Future Faculty Program}, Mississippi State University
              (2024--2025)
        \item \textbf{Certified Teaching Assistant, Level 2/3}, Mississippi State
              University (2024)
        \item \textbf{Fundamentals of Deep Learning}, Certificate of Competency,
              NVIDIA Deep Learning Institute (2023)
        \item \textbf{Psychological Counselling Program} (40 hours), Staff Development
              Centre, Uva Wellassa University / World Bank AHEAD Project (2021--2022)
        \item \textbf{Workshop: Writing Manuscripts and Proposals for Research Grants},
              AHEAD ELTA-ELSE Faculty Development Plan, Uva Wellassa University (2020)
        \item \textbf{Workshop: Online Teaching Tools and Methods for Effective
              Teaching and Assessment}, Uva Wellassa University (2020)
        \item \textbf{Certificate Course in Teaching in Higher Education},
              Staff Development Centre, Uva Wellassa University (2019)
    \end{highlights}
\end{onecolentry}
\vspace{0.18cm}

%--------------------------------------------------------------
\section{TECHNICAL SKILLS}
\vspace{0.08cm}
\begin{onecolentry}
    \textbf{Programming:} Python, C++, MATLAB \\
    \textbf{Deep learning:} PyTorch, TensorFlow/Keras, YOLO (v5--v26), SegFormer,
    instance and semantic segmentation, dataset design and annotation \\
    \textbf{Robotics and perception:} ROS, Gazebo, OpenCV, RGB-D and stereo
    vision (ZED 2i, RealSense), UR5e manipulator control, Clearpath Husky,
    GPS and map-based autonomous navigation, SLAM \\
    \textbf{Hardware and embedded:} Arduino, PLC, microcontroller systems,
    IoT sensor networks, instrumentation design
\end{onecolentry}
\vspace{0.18cm}

%--------------------------------------------------------------
\section{PROFESSIONAL MEMBERSHIPS}
\vspace{0.08cm}
\begin{onecolentry}
    \begin{highlights}
        \item Integrative Precision Agriculture Student Club (IPAC),
              University of Georgia, since 2025
        \item American Society of Agricultural and Biological Engineers (ASABE),
              Graduate Student Member, since 2024
        \item SPIE --- the International Society for Optics and Photonics,
              Student Member, since 2023
        \item Institute of Electrical and Electronics Engineers (IEEE),
              Graduate Student Member, since 2020
    \end{highlights}
\end{onecolentry}
\vspace{0.18cm}

%--------------------------------------------------------------
\section{RESEARCH SOFTWARE AND DATA}
\vspace{0.08cm}
\begin{onecolentry}
    \begin{highlights}
        \item \textbf{CottonID} (2026) --- web application for cotton boll detection,
              yield estimation, and defoliation-timing decision support
              (Next.js, FastAPI, RF-DETR/YOLO).
              % TODO: add a public link or repository if available
        \item \textbf{CottonBoll Harvest} (2026) --- supplementary code, trained-model
              resources, and supporting materials for the manuscript
              \emph{Selective Cotton Boll Localization for Robotic Harvesting:
              Evaluation of Deep Learning Vision Models under Field Conditions}.
              Repository private during manuscript preparation. Public release planned upon publication.
              \\ \repo{github.com/imtheva/CottonBoll\_Harvest}
        \item \textbf{CottonSim} (2025) --- open-source ROS/Gazebo simulator, virtual
              cotton farm, and vision-guided navigation stack for autonomous cotton
              picking. \\ \repo{github.com/imtheva/CottonSim}
        \item \textbf{Multi-Ripeness Blackberry} (2025) --- detection models and
              annotated field dataset (1{,}086 images, three ripeness classes, two
              seasons). \\ \repo{github.com/Zhanglab-abe/Multi-Ripeness\_Blackberry}
        \item \textbf{Catfish-Segment} (2023) --- semantic segmentation models for
              automated catfish processing.
              \\ \repo{github.com/Zhanglab-abe/Catfish-Segment}
    \end{highlights}
\end{onecolentry}
\vspace{0.18cm}

%--------------------------------------------------------------
\section{MEDIA COVERAGE}
\vspace{0.08cm}
\begin{onecolentry}
    \begin{highlights}
        \item \href{https://msstate-innovations.technologypublisher.com/technology/58899}{MSU
              Innovations}, ``Design and development of an autonomous vision-guided
              cotton picker for selective and targeted robotic picking''
              (October 2025)
        \item \href{https://farmflavor.com/mississippi/mississippi-technology/agriculture-technology-advancements-revolutionize-mississippi-ag-industry/}{Farm
              Flavor}, ``Agriculture technology advancements revolutionize
              Mississippi ag industry'' (November 2024)
        \item \href{https://www.wcbi.com/mississippi-state-researchers-develop-prototype-to-advance-cotton-industry/}{WCBI
              News}, ``Mississippi State researchers develop prototype to advance
              cotton industry'' (November 2023)
        \item \href{https://issuu.com/gradmarketing/docs/newsletter_-_ragged/40}{MSU
              Graduate School Newsletter}, award announcement (October 2023)
    \end{highlights}
\end{onecolentry}
\vspace{0.18cm}

%--------------------------------------------------------------
% \section{REFERENCES}
% -------------------------------------------------------------------------
% REFERENCES SECTION -- COMMENTED OUT.
% Uncomment and fill in before submitting. For faculty applications, naming
% 3-5 referees with full contact details is the convention.
%
% \vspace{0.08cm}
% \begin{onecolentry}
% \begin{tabular}{@{}p{8.2cm} p{8.2cm}@{}}
% \textbf{Dr.\ Xin Zhang} & \textbf{Name} \\
% Ph.D.\ Advisor; Rank & Title \\
% School of Environmental, Civil, Agricultural & Department \\
% \& Mechanical Engineering & Institution \\
% University of Georgia & email \\
% email \textbar{} phone & phone \\
% \end{tabular}
% \end{onecolentry}

\end{document}
